% year: תשפ"ב
% type: מבחן
% semester: א
% moed: מיוחד
% number: 6ב
% points: 9
% categories: improper-integrals
% source: exams/מבחנים/תשפב/מועד ג תשפב.pdf

\begin{question}
בדקו התכנסות של האינטגרל הלא אמיתי
\[ \int_1^\infty\frac{e^x}{e^{2x}+1}dx \]
\end{question}

\begin{hint}
הפונקציה הקדומה היא $\arctan(e^x)$ (הצבה $t=e^x$).
\end{hint}

\begin{solution}
נציב $t=e^x$, $dt=e^xdx$: $\int\frac{e^x}{e^{2x}+1}dx=\int\frac{dt}{t^2+1}=\arctan(e^x)+C$. לכן
\[ \int_1^R\frac{e^x}{e^{2x}+1}dx=\arctan(e^R)-\arctan(e). \]
כאשר $R\to\infty$: $e^R\to\infty$ ו-$\lim_{t\to\infty}\arctan t=\frac\pi2$, ולכן
\[ \int_1^\infty\frac{e^x}{e^{2x}+1}dx=\lim_{R\to\infty}\left(\arctan(e^R)-\arctan e\right)=\boxed{\frac\pi2-\arctan e}\approx0.3466. \]
האינטגרל \textbf{מתכנס}.

(לחלופין, ללא חישוב: $0<\frac{e^x}{e^{2x}+1}<\frac{e^x}{e^{2x}}=e^{-x}$, ו-$\int_1^\infty e^{-x}dx=e^{-1}$ מתכנס, ולכן לפי מבחן ההשוואה האינטגרל מתכנס.)
\end{solution}
